\documentclass[11pt]{article}
\usepackage{icmlko}
\begin{document}

\title{판이 두 개였다 --- 2년간 우리가 순위라 부른 것은 두 수였고, 채택 넷 중 둘은 어느 쪽을 쓰느냐로 부호가 갈렸다}
\author{월드모델 실험실}
\date{노트 526 · 2026년 8월 4일}
\maketitle

\begin{abstract}
용어 대응표를 쓰다가 \texttt{lab/harness.py} 의 \texttt{board()} 가
\textbf{도메인 단순 평균}인데, 사이클 스크립트가 2년째 ``판''이라 불러 온 것은
\textbf{유보 행수 가중 평균}임을 발견했다. 같은 이름, 다른 수다.

\begin{center}
\begin{tabular}{lrrrr}
\hline
채택 & 가중 판 & 부호 & 단순 판 & 부호 \\
\hline
\texttt{DOMDROP}(498) & $\mathbf{+0.0043}$ & $12/12$ & $\mathbf{-0.0007}$ & $5/12$ \\
\texttt{TRAINCAP}(509) & $+0.0018$ & $8/12$ & $+0.0001$ & $5/12$ \\
\texttt{TRAINW}(514) & $+0.0005$ & $7/12$ & $+0.0012$ & $9/12$ \\
\texttt{TRAINR}(516) & $\mathbf{+0.0028}$ & $10/12$ & $\mathbf{-0.0010}$ & $4/12$ \\
\hline
\end{tabular}
\end{center}

\textbf{넷 중 둘이 부호가 뒤집힌다.} 챔피언 자체도 가중 $+0.4908$ 대 단순
$+0.4507$ 로 $\mathbf{0.0401}$ 벌어지는데, 이는 \textbf{역대 채택 22건 중
가장 큰 것($+0.0088$)의 4.6배}다. 어느 정의를 쓰느냐가 2년간 번 것 전부보다
크다. 그래서 취향이 아니라 \textbf{목표}로 갈랐다 --- 제품 목표에 가장 가까운
자, 곧 \textbf{집 밖 짝}으로.
\end{abstract}

\begin{figure}[t]
\centering
\includegraphics[width=\textwidth]{figs/f1.pdf}
\caption{\textbf{가} --- 같은 채택, 다른 정의. \textbf{나} --- 집 밖 짝과의
부호 일치. \textbf{다} --- 왜, 그리고 한계.}
\end{figure}

\section{집 밖 짝이 가른다}

같은 실행 $\cdot$ 같은 씨앗(짝 $12$)에서 채택 넷을 되돌리며 여섯 자를 함께
쟀다. 집 밖 짝은 학습에 한 줄도 없는 수집이라 \emph{도메인 무게와 무관}하고,
``처음 보는 시장으로 건너간다''는 목표에 가장 가깝다.

\begin{center}
\begin{tabular}{lrrl}
\hline
채택 & KR 만화 & 비게임 앱 & 어느 정의와 부호가 맞나 \\
\hline
\texttt{DOMDROP} & $+0.0016$ & $+0.0040$ & \textbf{가중만} \\
\texttt{TRAINCAP} & $+0.0032$ & $+0.0058$ & 둘 다 \\
\texttt{TRAINW} & $+0.0159$ & $+0.0066$ & 둘 다 \\
\texttt{TRAINR} & $+0.0159$ & $+0.0077$ & \textbf{가중만} \\
\hline
\end{tabular}
\end{center}

\begin{quote}
\textbf{가중 $8/8$ $\cdot$ 단순 $4/8$.} 그리고 순위 상관 대신 \emph{쌍 정확도}로
바꿔도 똑같다(쌍가중 $8$ $\cdot$ 쌍단순 $4$) --- \textbf{가르는 것은 통계량이
아니라 무게다.}
\end{quote}

\section{왜 그런가 --- 첫 양수를 그대로 받지 않는다}

세 갈래로 갈랐다.

\paragraph{① 가중은 역분산 가중에 가깝다.} 스피어만 $\rho$ 의 표준오차는
대략 $(1-\rho^2)/\sqrt{n-1}$ 이다. 이걸로 최소분산 무게를 만들어 행수 가중과
견주면 \textbf{순위상관 $+0.936$}. 즉 ``도메인마다 같은 효과가 있다''를 재는
상황에서 행수 가중은 \emph{거의 최적}이다.

\paragraph{② 판 자체가 덜 시끄럽다.} 판의 표준오차는 가중 $0.0129$ 대 단순
$0.0197$ --- \textbf{단순이 $1.53$ 배 시끄럽다.} 단순 평균은 팝업($65$행,
표준오차 $0.115$)에 웹툰($711$행, $0.029$)의 $4.8$ 배 정밀도를 가진 것처럼
무게를 준다.

\paragraph{③ 단순은 ``빌려 쓰는'' 도메인에 무게를 몰아준다.} 노트 505 는
도메인마다 \emph{자기 천장}(그 도메인 자기 행만으로 교차검증)을 쟀다. 자기
천장의 $1.7$ 배 이상을 받거나 천장이 음수인 도메인 다섯(팝업 $\cdot$ 아이돌
$\cdot$ 펀딩 $\cdot$ 만화 $\cdot$ 시장팝업)은 점수의 대부분이 \textbf{남에게서
온 것}이다.

\begin{center}
\begin{tabular}{lrr}
\hline
 & 가중 & 단순 \\
\hline
빌려 쓰는 다섯의 무게 & $30.0\%$ & $\mathbf{45.5\%}$ ($1.52$ 배) \\
\hline
\end{tabular}
\end{center}

그 도메인을 올리는 변경은 \emph{그 도메인 안에서만 통하는 것}을 좇게 되고,
그래서 전이와 어긋난다.

\begin{quote}
\textbf{정직하게} --- ``작은 도메인이 빌린다''는 아니다. 유보 행수와 빌린
배율의 순위상관은 $-0.433$ 에 그치고 펀딩($529$행 $\cdot$ $1.72$ 배)이 반례다.
\textbf{작아서가 아니라 빌려서}라고 써야 맞다.
\end{quote}

\section{그래도 이것은 목표 상대적이다}

가중이 이긴 것은 \textbf{우리 목표가 전이이기 때문}이다. 두 정의는 같은 것을
다르게 재는 것이 아니라 \textbf{다른 것을 잰다} --- 가중은 ``무작위로 고른 한
레코드''를, 단순은 ``무작위로 고른 한 도메인''을 잰다.

\begin{quote}
\textbf{조항} --- \textbf{두 지표가 갈리면 더 나은 추정량을 고르는 문제가
아니라 무엇을 재고 싶은지를 정하는 문제다.} 우리 목표는 도메인을 가로지르는
구조이므로 \textbf{순위는 유보 행수 가중 판으로 고정한다}. 도메인을 똑같이
아끼는 것이 목표인 조직이라면 단순이 옳고, 그 조직은 우리 채택 넷 중 둘을
안 가져가야 한다.
\end{quote}

\section*{배선 하나를 고쳤다}

첫 측정(\texttt{n636})은 되돌림 속성을 \emph{예보 전에} 복원했다.
\texttt{TRAINCAP} $\cdot$ \texttt{TRAINW} $\cdot$ \texttt{TRAINR} 은 학습
전용이라 무관하지만 \texttt{DOMDROP} 은 적합과 예보 양쪽에 걸리므로, 그 팔은
``가림 없이 학습하고 가림 쓰고 예보''라는 어긋난 배선이었다. 고친 뒤
$+0.0022 \to +0.0043$ 이다. \textbf{부호가 갈린다는 결론은 그대로다.}
(자기 조항 --- 학습 쪽 변경은 학습 자료를 흔들어 확인한다 --- 을 또 밟았다.
이번엔 다음 노트가 잡았다.)

\section*{남은 것}

\begin{itemize}
\item \textbf{원장의 판 수치에 정의를 명시한다.} 지금까지 기록은 전부 가중인데
그렇게 안 적혀 있다.
\item \textbf{\texttt{harness.board()} 를 고치지 않는다} --- 배포 파이프라인이
쓰는 수라 손대면 다른 것이 움직인다. 대신 \textbf{이름을 갈라} 적는다.
\item \textbf{도메인 무게가 자라면 다시 본다} --- 유보가 늘면 무게가 바뀌고,
그러면 이 판정도 다시 재야 한다(노트 406 계열 조항).
\end{itemize}

\end{document}
